\documentclass[11pt,letterpaper]{article}

% Paquetes necesarios
\usepackage[utf8]{inputenc}
\usepackage[spanish]{babel}
\usepackage[margin=2cm]{geometry}
\usepackage{fancyhdr}
\usepackage{graphicx}
\usepackage{listings}
\usepackage{xcolor}
\usepackage{hyperref}
\usepackage{tabularx}
\usepackage{array}

% Configuración del header y footer
\pagestyle{fancy}
\fancyhf{}
\lhead{NLPY-1: Fundamentos de Python}
\rhead{Semana 8: Programación Orientada a Objetos}
\rfoot{Página \thepage}

% Configuración de listings para Python
\lstset{
    language=Python,
    basicstyle=\ttfamily\small,
    keywordstyle=\color{blue}\bfseries,
    stringstyle=\color{red},
    commentstyle=\color{gray}\itshape,
    numbers=left,
    numberstyle=\tiny\color{gray},
    stepnumber=1,
    numbersep=8pt,
    showstringspaces=false,
    breaklines=true,
    frame=single,
    backgroundcolor=\color{gray!10},
    tabsize=4,
    captionpos=b
}

% Configuración de hyperref
\hypersetup{
    colorlinks=true,
    linkcolor=blue,
    urlcolor=blue,
    citecolor=blue
}

\begin{document}

% Título
\begin{center}
    {\LARGE \textbf{Tarea 8: Sistema de Gestión de Biblioteca}}\\[0.3cm]
    {\Large Semana 8: Programación Orientada a Objetos}\\[0.2cm]
    {\large NLPY-1: Fundamentos de Python}\\[0.2cm]
    {\large nablabs}\\[0.5cm]
\end{center}

\section{Introducción}

La \textbf{Programación Orientada a Objetos (POO)} es uno de los paradigmas más importantes y utilizados en el desarrollo de software moderno. Permite modelar el mundo real mediante objetos que contienen datos (atributos) y comportamientos (métodos), facilitando la organización, reutilización y mantenimiento del código.

Python es un lenguaje multiparadigma que soporta POO de manera elegante y poderosa. En esta semana aprenderás a diseñar sistemas complejos usando clases, herencia, encapsulamiento, polimorfismo y otros conceptos fundamentales de la POO.

\subsection{¿Por qué Programación Orientada a Objetos?}

\begin{itemize}
    \item \textbf{Modelado natural}: Representa entidades del mundo real como objetos
    \item \textbf{Reutilización de código}: Herencia y composición evitan duplicación
    \item \textbf{Encapsulamiento}: Protege datos y controla el acceso a ellos
    \item \textbf{Modularidad}: Código organizado en unidades lógicas independientes
    \item \textbf{Mantenibilidad}: Cambios localizados, menos propenso a errores
    \item \textbf{Polimorfismo}: Mismo comportamiento con diferentes implementaciones
\end{itemize}

\subsection{Objetivos de Aprendizaje}

Al completar esta tarea, el estudiante será capaz de:

\begin{itemize}
    \item Diseñar e implementar clases con atributos y métodos
    \item Aplicar el método constructor \texttt{\_\_init\_\_()}
    \item Implementar métodos especiales (\texttt{\_\_str\_\_}, \texttt{\_\_repr\_\_}, etc.)
    \item Utilizar propiedades con decoradores \texttt{@property} y setters
    \item Aplicar herencia para reutilizar y extender funcionalidad
    \item Implementar polimorfismo mediante herencia
    \item Usar métodos de clase (\texttt{@classmethod}) y estáticos (\texttt{@staticmethod})
    \item Aplicar encapsulamiento con atributos privados y protegidos
    \item Comprender y usar composición de objetos
    \item Diseñar jerarquías de clases coherentes
\end{itemize}

\section{Descripción del Proyecto}

Desarrollarás un \textbf{Sistema de Gestión de Biblioteca} completo que permita administrar libros, usuarios, préstamos y multas usando POO. El sistema debe implementar múltiples clases relacionadas con herencia, composición y polimorfismo.

\subsection{Funcionalidades Requeridas}

\subsubsection{1. Sistema de Clases (40 puntos - Funcionalidad)}

Tu sistema debe implementar la siguiente jerarquía de clases:

\textbf{A) Clase Base: Persona (5 puntos)}
\begin{itemize}
    \item Atributos: nombre, identificación, email, teléfono
    \item Métodos: \texttt{\_\_init\_\_}, \texttt{\_\_str\_\_}, \texttt{\_\_repr\_\_}
    \item Validación de datos en el constructor
    \item Método para actualizar información de contacto
\end{itemize}

\textbf{B) Clases Derivadas de Persona (8 puntos)}
\begin{itemize}
    \item \textbf{Usuario}: estudiante o miembro regular
    \begin{itemize}
        \item Atributos adicionales: tipo\_usuario, libros\_prestados, historial
        \item Límite de préstamos según tipo (estudiante: 3, regular: 5)
        \item Método para verificar si puede prestar más libros
    \end{itemize}
    \item \textbf{Bibliotecario}: empleado de la biblioteca
    \begin{itemize}
        \item Atributos: número\_empleado, turno, especialidad
        \item Métodos para gestionar préstamos y multas
    \end{itemize}
\end{itemize}

\textbf{C) Clase Libro (8 puntos)}
\begin{itemize}
    \item Atributos: ISBN, título, autor, editorial, año, categoría
    \item Atributos de estado: disponible, prestado\_a, fecha\_prestamo
    \item Métodos: prestar, devolver, obtener\_info
    \item Propiedades: \texttt{@property} para estado y días\_prestado
    \item Validación de ISBN
\end{itemize}

\textbf{D) Clase Prestamo (7 puntos)}
\begin{itemize}
    \item Atributos: usuario, libro, fecha\_prestamo, fecha\_devolucion, devuelto
    \item Calcular días de préstamo y días de retraso
    \item Método para registrar devolución
    \item Calcular multa por retraso (₡500 por día)
    \item Usar composición (tiene Usuario y Libro)
\end{itemize}

\textbf{E) Clase Biblioteca (12 puntos)}
\begin{itemize}
    \item Gestiona colecciones de libros, usuarios y préstamos
    \item Métodos CRUD para libros y usuarios
    \item Registrar préstamos y devoluciones
    \item Buscar libros por título, autor, ISBN o categoría
    \item Generar reportes y estadísticas
    \item Listar libros disponibles/prestados
    \item Calcular total de multas pendientes
\end{itemize}

\subsubsection{2. Requisitos Técnicos}

\textbf{Uso de POO:}
\begin{itemize}
    \item Implementar herencia correctamente
    \item Usar \texttt{super()} para llamar constructores padre
    \item Implementar al menos 3 métodos especiales (\texttt{\_\_str\_\_}, \texttt{\_\_repr\_\_}, \texttt{\_\_eq\_\_})
    \item Usar \texttt{@property} y setters para atributos calculados
    \item Implementar al menos un \texttt{@classmethod} y un \texttt{@staticmethod}
\end{itemize}

\textbf{Encapsulamiento:}
\begin{itemize}
    \item Atributos privados (\_\_) para datos sensibles
    \item Atributos protegidos (\_) para datos internos
    \item Getters y setters con validación
    \item No acceder directamente a atributos desde fuera
\end{itemize}

\textbf{Diseño:}
\begin{itemize}
    \item Jerarquía de clases coherente y lógica
    \item Responsabilidad única por clase
    \item Composición sobre herencia donde sea apropiado
    \item Métodos bien nombrados y documentados
\end{itemize}

\section{Conceptos de POO en Python}

\subsection{Elementos Fundamentales}

\textbf{Clases y Objetos:}
\begin{lstlisting}[caption={Definición básica de clase}]
class Libro:
    # Constructor
    def __init__(self, titulo, autor, isbn):
        self.titulo = titulo
        self.autor = autor
        self.isbn = isbn
        self.disponible = True
    
    # Metodo de instancia
    def prestar(self):
        if self.disponible:
            self.disponible = False
            return True
        return False
    
    # Metodo especial para representacion
    def __str__(self):
        return f"{self.titulo} por {self.autor}"

# Crear objetos (instancias)
libro1 = Libro("1984", "George Orwell", "978-0451524935")
print(libro1)  # Usa __str__
\end{lstlisting}

\textbf{Herencia:}
\begin{lstlisting}[caption={Herencia de clases}]
class Persona:
    def __init__(self, nombre, identificacion):
        self.nombre = nombre
        self.identificacion = identificacion
    
    def presentarse(self):
        return f"Hola, soy {self.nombre}"

class Usuario(Persona):  # Usuario hereda de Persona
    def __init__(self, nombre, identificacion, tipo):
        super().__init__(nombre, identificacion)  # Llamar constructor padre
        self.tipo = tipo
        self.libros_prestados = []
    
    def presentarse(self):  # Sobreescritura (override)
        return f"{super().presentarse()}, usuario {self.tipo}"
\end{lstlisting}

\textbf{Propiedades:}
\begin{lstlisting}[caption={Uso de @property}]
class Libro:
    def __init__(self, titulo, precio):
        self.titulo = titulo
        self._precio = precio  # Atributo protegido
    
    @property
    def precio(self):
        """Getter para precio"""
        return self._precio
    
    @precio.setter
    def precio(self, valor):
        """Setter con validacion"""
        if valor < 0:
            raise ValueError("Precio no puede ser negativo")
        self._precio = valor
    
    @property
    def precio_con_iva(self):
        """Propiedad calculada (solo lectura)"""
        return self._precio * 1.13

# Uso
libro = Libro("Python 101", 15000)
print(libro.precio)  # Usa getter
libro.precio = 18000  # Usa setter
print(libro.precio_con_iva)  # Propiedad calculada
\end{lstlisting}

\section{Plantilla de Código}

\subsection{Estructura Base del Sistema}

\begin{lstlisting}[caption={Plantilla de clases principales}]
from datetime import datetime, timedelta
from typing import List, Optional

# ==================================================================
# CLASE BASE: PERSONA
# ==================================================================

class Persona:
    """Clase base para representar una persona."""
    
    def __init__(self, nombre: str, identificacion: str, 
                 email: str, telefono: str):
        """
        Constructor de Persona.
        
        Args:
            nombre: Nombre completo
            identificacion: Cedula o ID
            email: Correo electronico
            telefono: Numero de telefono
        """
        # TODO: Validar que los datos no esten vacios
        # TODO: Asignar atributos
        pass
    
    def __str__(self) -> str:
        """Representacion en string."""
        # TODO: Retornar representacion legible
        pass
    
    def __repr__(self) -> str:
        """Representacion tecnica."""
        # TODO: Retornar representacion para debugging
        pass
    
    def actualizar_contacto(self, email: str = None, 
                           telefono: str = None):
        """Actualiza informacion de contacto."""
        # TODO: Actualizar email y/o telefono si se proveen
        pass

# ==================================================================
# CLASES DERIVADAS DE PERSONA
# ==================================================================

class Usuario(Persona):
    """Usuario de la biblioteca."""
    
    LIMITE_PRESTAMOS = {
        'estudiante': 3,
        'regular': 5
    }
    
    def __init__(self, nombre: str, identificacion: str, 
                 email: str, telefono: str, tipo_usuario: str):
        """
        Constructor de Usuario.
        
        Args:
            tipo_usuario: 'estudiante' o 'regular'
        """
        # TODO: Llamar constructor de Persona con super()
        # TODO: Validar tipo_usuario
        # TODO: Inicializar libros_prestados como lista vacia
        # TODO: Inicializar historial como lista vacia
        pass
    
    @property
    def puede_prestar(self) -> bool:
        """Verifica si el usuario puede prestar mas libros."""
        # TODO: Comparar cantidad de libros con el limite
        pass
    
    @property
    def cantidad_prestamos(self) -> int:
        """Retorna cantidad de libros actualmente prestados."""
        # TODO: Retornar longitud de libros_prestados
        pass
    
    def agregar_prestamo(self, libro):
        """Agrega un libro a los prestados."""
        # TODO: Agregar a libros_prestados
        pass
    
    def remover_prestamo(self, libro):
        """Remueve un libro de los prestados."""
        # TODO: Remover de libros_prestados
        pass

class Bibliotecario(Persona):
    """Empleado bibliotecario."""
    
    def __init__(self, nombre: str, identificacion: str, 
                 email: str, telefono: str, 
                 numero_empleado: str, turno: str):
        """Constructor de Bibliotecario."""
        # TODO: Llamar constructor de Persona
        # TODO: Asignar numero_empleado y turno
        pass
    
    def registrar_prestamo(self, prestamo) -> bool:
        """Registra un prestamo en el sistema."""
        # TODO: Validar y registrar prestamo
        pass
    
    def registrar_devolucion(self, prestamo) -> dict:
        """Registra una devolucion y calcula multa."""
        # TODO: Procesar devolucion
        # TODO: Calcular multa si hay retraso
        pass

# ==================================================================
# CLASE LIBRO
# ==================================================================

class Libro:
    """Representa un libro en la biblioteca."""
    
    def __init__(self, isbn: str, titulo: str, autor: str,
                 editorial: str, año: int, categoria: str):
        """Constructor de Libro."""
        # TODO: Validar ISBN (formato basico)
        # TODO: Asignar todos los atributos
        # TODO: Inicializar disponible = True
        # TODO: Inicializar prestado_a = None
        # TODO: Inicializar fecha_prestamo = None
        pass
    
    @property
    def disponible(self) -> bool:
        """Retorna si el libro esta disponible."""
        # TODO: Retornar estado de disponibilidad
        pass
    
    @property
    def dias_prestado(self) -> int:
        """Calcula dias que lleva prestado el libro."""
        # TODO: Si esta prestado, calcular dias desde fecha_prestamo
        # TODO: Si no, retornar 0
        pass
    
    def prestar(self, usuario) -> bool:
        """Presta el libro a un usuario."""
        # TODO: Verificar que este disponible
        # TODO: Marcar como no disponible
        # TODO: Asignar usuario y fecha
        pass
    
    def devolver(self) -> bool:
        """Marca el libro como devuelto."""
        # TODO: Marcar como disponible
        # TODO: Limpiar prestado_a y fecha_prestamo
        pass
    
    def __str__(self) -> str:
        """Representacion en string."""
        # TODO: Retornar info del libro
        pass
    
    def __eq__(self, otro) -> bool:
        """Compara libros por ISBN."""
        # TODO: Comparar ISBN de ambos libros
        pass
    
    @staticmethod
    def validar_isbn(isbn: str) -> bool:
        """Valida formato basico de ISBN."""
        # TODO: Verificar longitud y formato
        pass

# ==================================================================
# CLASE PRESTAMO
# ==================================================================

class Prestamo:
    """Representa un prestamo de libro."""
    
    DIAS_LIMITE = 14  # 2 semanas
    MULTA_POR_DIA = 500  # Colones
    
    def __init__(self, usuario: Usuario, libro: Libro):
        """Constructor de Prestamo."""
        # TODO: Asignar usuario y libro
        # TODO: Asignar fecha_prestamo (datetime.now())
        # TODO: Calcular fecha_devolucion (fecha + DIAS_LIMITE)
        # TODO: Inicializar devuelto = False
        # TODO: Inicializar fecha_devolucion_real = None
        pass
    
    @property
    def dias_transcurridos(self) -> int:
        """Calcula dias desde el prestamo."""
        # TODO: Calcular diferencia entre hoy y fecha_prestamo
        pass
    
    @property
    def dias_retraso(self) -> int:
        """Calcula dias de retraso."""
        # TODO: Si no esta devuelto, calcular retraso desde fecha limite
        # TODO: Si esta devuelto, calcular desde fecha real
        # TODO: Retornar 0 si no hay retraso
        pass
    
    @property
    def multa(self) -> float:
        """Calcula la multa por retraso."""
        # TODO: Multiplicar dias_retraso por MULTA_POR_DIA
        pass
    
    def registrar_devolucion(self) -> dict:
        """Registra la devolucion del libro."""
        # TODO: Marcar como devuelto
        # TODO: Asignar fecha_devolucion_real
        # TODO: Retornar dict con info de multa
        pass
    
    def __str__(self) -> str:
        """Representacion en string."""
        # TODO: Retornar info del prestamo
        pass

# ==================================================================
# CLASE BIBLIOTECA
# ==================================================================

class Biblioteca:
    """Sistema de gestion de biblioteca."""
    
    def __init__(self, nombre: str):
        """Constructor de Biblioteca."""
        # TODO: Asignar nombre
        # TODO: Inicializar listas: libros, usuarios, prestamos
        pass
    
    # --- Gestion de Libros ---
    
    def agregar_libro(self, libro: Libro) -> bool:
        """Agrega un libro al catalogo."""
        # TODO: Verificar que ISBN no exista
        # TODO: Agregar a lista de libros
        pass
    
    def buscar_libro_por_isbn(self, isbn: str) -> Optional[Libro]:
        """Busca un libro por ISBN."""
        # TODO: Buscar en lista de libros
        pass
    
    def buscar_libros_por_titulo(self, titulo: str) -> List[Libro]:
        """Busca libros por titulo (coincidencia parcial)."""
        # TODO: Filtrar libros que contengan el titulo
        pass
    
    def buscar_libros_por_autor(self, autor: str) -> List[Libro]:
        """Busca libros por autor."""
        # TODO: Filtrar libros del autor
        pass
    
    def listar_libros_disponibles(self) -> List[Libro]:
        """Lista todos los libros disponibles."""
        # TODO: Filtrar libros con disponible=True
        pass
    
    # --- Gestion de Usuarios ---
    
    def registrar_usuario(self, usuario: Usuario) -> bool:
        """Registra un nuevo usuario."""
        # TODO: Verificar que identificacion no exista
        # TODO: Agregar a lista de usuarios
        pass
    
    def buscar_usuario(self, identificacion: str) -> Optional[Usuario]:
        """Busca un usuario por identificacion."""
        # TODO: Buscar en lista de usuarios
        pass
    
    # --- Gestion de Prestamos ---
    
    def realizar_prestamo(self, identificacion_usuario: str, 
                         isbn: str) -> dict:
        """Realiza un prestamo de libro."""
        # TODO: Buscar usuario y libro
        # TODO: Verificar que usuario puede prestar
        # TODO: Verificar que libro este disponible
        # TODO: Crear objeto Prestamo
        # TODO: Actualizar estado de libro y usuario
        # TODO: Agregar a lista de prestamos
        pass
    
    def realizar_devolucion(self, isbn: str) -> dict:
        """Procesa la devolucion de un libro."""
        # TODO: Buscar prestamo activo del libro
        # TODO: Registrar devolucion
        # TODO: Actualizar libro y usuario
        # TODO: Retornar info con multa
        pass
    
    # --- Reportes y Estadisticas ---
    
    def generar_reporte_general(self) -> dict:
        """Genera reporte general de la biblioteca."""
        # TODO: Contar total de libros, usuarios, prestamos
        # TODO: Calcular libros disponibles/prestados
        # TODO: Calcular total de multas pendientes
        pass
    
    def libros_mas_prestados(self, limite: int = 5) -> List[tuple]:
        """Retorna los libros mas prestados."""
        # TODO: Contar prestamos por libro
        # TODO: Ordenar y retornar top N
        pass
    
    @classmethod
    def crear_biblioteca_ejemplo(cls):
        """Crea una biblioteca con datos de ejemplo."""
        # TODO: Crear instancia de Biblioteca
        # TODO: Agregar libros y usuarios de ejemplo
        # TODO: Retornar biblioteca poblada
        pass
    
    def __str__(self) -> str:
        """Representacion en string."""
        # TODO: Retornar info basica de la biblioteca
        pass

# ==================================================================
# MENU INTERACTIVO
# ==================================================================

def menu_principal():
    """Menu principal del sistema."""
    # TODO: Crear instancia de Biblioteca
    # TODO: Implementar menu con opciones:
    #   1. Gestion de libros
    #   2. Gestion de usuarios
    #   3. Realizar prestamo
    #   4. Realizar devolucion
    #   5. Buscar libros
    #   6. Reportes
    #   7. Salir
    pass

# ==================================================================
# MAIN
# ==================================================================

if __name__ == "__main__":
    print("Sistema de Gestion de Biblioteca")
    print("NLPY-1 - Semana 8: Programacion Orientada a Objetos")
    menu_principal()
\end{lstlisting}

\subsection{Ejemplo: Implementación de la Clase Libro}

Como referencia, aquí está la implementación completa de la clase Libro:

\begin{lstlisting}[caption={Ejemplo: Clase Libro completa}]
from datetime import datetime

class Libro:
    """Representa un libro en la biblioteca."""
    
    def __init__(self, isbn: str, titulo: str, autor: str,
                 editorial: str, año: int, categoria: str):
        """Constructor de Libro."""
        if not self.validar_isbn(isbn):
            raise ValueError(f"ISBN invalido: {isbn}")
        
        self.isbn = isbn
        self.titulo = titulo
        self.autor = autor
        self.editorial = editorial
        self.año = año
        self.categoria = categoria
        self._disponible = True
        self._prestado_a = None
        self._fecha_prestamo = None
    
    @property
    def disponible(self) -> bool:
        """Retorna si el libro esta disponible."""
        return self._disponible
    
    @property
    def dias_prestado(self) -> int:
        """Calcula dias que lleva prestado el libro."""
        if not self._disponible and self._fecha_prestamo:
            delta = datetime.now() - self._fecha_prestamo
            return delta.days
        return 0
    
    def prestar(self, usuario) -> bool:
        """Presta el libro a un usuario."""
        if not self._disponible:
            return False
        
        self._disponible = False
        self._prestado_a = usuario
        self._fecha_prestamo = datetime.now()
        return True
    
    def devolver(self) -> bool:
        """Marca el libro como devuelto."""
        if self._disponible:
            return False
        
        self._disponible = True
        self._prestado_a = None
        self._fecha_prestamo = None
        return True
    
    def __str__(self) -> str:
        """Representacion en string."""
        estado = "Disponible" if self._disponible else "Prestado"
        return f"'{self.titulo}' por {self.autor} ({self.año}) - {estado}"
    
    def __repr__(self) -> str:
        """Representacion tecnica."""
        return f"Libro(isbn='{self.isbn}', titulo='{self.titulo}')"
    
    def __eq__(self, otro) -> bool:
        """Compara libros por ISBN."""
        if not isinstance(otro, Libro):
            return False
        return self.isbn == otro.isbn
    
    @staticmethod
    def validar_isbn(isbn: str) -> bool:
        """Valida formato basico de ISBN."""
        # ISBN-10: 10 digitos
        # ISBN-13: 13 digitos (con guiones opcionales)
        isbn_limpio = isbn.replace('-', '').replace(' ', '')
        return len(isbn_limpio) in [10, 13] and isbn_limpio.isdigit()

# Ejemplo de uso:
libro = Libro(
    isbn="978-0451524935",
    titulo="1984",
    autor="George Orwell",
    editorial="Signet Classic",
    año=1950,
    categoria="Ficcion"
)

print(libro)  # '1984' por George Orwell (1950) - Disponible
print(repr(libro))  # Libro(isbn='978-0451524935', titulo='1984')
\end{lstlisting}

\section{Ejemplos de Ejecución}

\subsection{Creación de Objetos y Herencia}

\begin{lstlisting}[caption={Ejemplo de creación de usuarios y bibliotecarios}]
# Crear usuario estudiante
estudiante = Usuario(
    nombre="Maria Rodriguez",
    identificacion="1-2345-6789",
    email="maria.rodriguez@estudiante.cr",
    telefono="8888-8888",
    tipo_usuario="estudiante"
)

print(estudiante)
# Salida: Maria Rodriguez (1-2345-6789) - Usuario estudiante

# Verificar limite de prestamos
print(f"Puede prestar: {estudiante.puede_prestar}")  # True
print(f"Prestamos actuales: {estudiante.cantidad_prestamos}")  # 0

# Crear bibliotecario
bibliotecario = Bibliotecario(
    nombre="Carlos Mora",
    identificacion="3-456-789012",
    email="cmora@biblioteca.cr",
    telefono="2222-2222",
    numero_empleado="BIB-001",
    turno="matutino"
)

print(bibliotecario)
# Salida: Carlos Mora (3-456-789012) - Bibliotecario BIB-001
\end{lstlisting}

\subsection{Gestión de Libros y Préstamos}

\begin{lstlisting}[caption={Ejemplo de operaciones con libros}]
# Crear biblioteca
biblioteca = Biblioteca("Biblioteca Central nablabs")

# Agregar libros
libro1 = Libro(
    isbn="978-0451524935",
    titulo="1984",
    autor="George Orwell",
    editorial="Signet Classic",
    año=1950,
    categoria="Ficcion"
)

libro2 = Libro(
    isbn="978-0062315007",
    titulo="Sapiens",
    autor="Yuval Noah Harari",
    editorial="Harper",
    año=2015,
    categoria="Historia"
)

biblioteca.agregar_libro(libro1)
biblioteca.agregar_libro(libro2)

# Registrar usuario
biblioteca.registrar_usuario(estudiante)

# Realizar prestamo
resultado = biblioteca.realizar_prestamo(
    identificacion_usuario="1-2345-6789",
    isbn="978-0451524935"
)

print(resultado)
# Salida:
# {
#     'exito': True,
#     'mensaje': 'Prestamo realizado exitosamente',
#     'libro': '1984',
#     'usuario': 'Maria Rodriguez',
#     'fecha_devolucion': '2024-11-07',
#     'dias_limite': 14
# }

# Verificar estado
print(f"Libro disponible: {libro1.disponible}")  # False
print(f"Dias prestado: {libro1.dias_prestado}")  # 0 (recien prestado)
print(f"Usuario puede prestar: {estudiante.puede_prestar}")  # True (2/3)
\end{lstlisting}

\subsection{Devolución y Cálculo de Multas}

\begin{lstlisting}[caption={Ejemplo de devolución con multa}]
# Simular que pasaron 20 dias (14 dias limite + 6 de retraso)
# En la practica esto seria automatico por las fechas

resultado = biblioteca.realizar_devolucion(isbn="978-0451524935")

print(resultado)
# Salida:
# {
#     'exito': True,
#     'mensaje': 'Devolucion registrada',
#     'libro': '1984',
#     'dias_prestamo': 20,
#     'dias_retraso': 6,
#     'multa': 3000,  # 6 dias * ₡500
#     'tiene_multa': True
# }

# Libro vuelve a estar disponible
print(f"Libro disponible: {libro1.disponible}")  # True
\end{lstlisting}

\subsection{Búsquedas y Consultas}

\begin{lstlisting}[caption={Ejemplo de búsquedas}]
# Buscar libros por titulo
resultados = biblioteca.buscar_libros_por_titulo("sapiens")
for libro in resultados:
    print(libro)
# Salida: 'Sapiens' por Yuval Noah Harari (2015) - Disponible

# Buscar por autor
resultados = biblioteca.buscar_libros_por_autor("Orwell")
print(f"Libros de Orwell: {len(resultados)}")
# Salida: Libros de Orwell: 1

# Listar disponibles
disponibles = biblioteca.listar_libros_disponibles()
print(f"\nLibros disponibles: {len(disponibles)}")
for libro in disponibles:
    print(f"  - {libro.titulo}")
# Salida:
# Libros disponibles: 2
#   - 1984
#   - Sapiens
\end{lstlisting}

\subsection{Reportes y Estadísticas}

\begin{lstlisting}[caption={Ejemplo de generación de reportes}]
# Generar reporte general
reporte = biblioteca.generar_reporte_general()

print("\n" + "="*60)
print(f"REPORTE DE BIBLIOTECA: {biblioteca.nombre}")
print("="*60)
print(f"Total de libros: {reporte['total_libros']}")
print(f"Libros disponibles: {reporte['libros_disponibles']}")
print(f"Libros prestados: {reporte['libros_prestados']}")
print(f"Total de usuarios: {reporte['total_usuarios']}")
print(f"Prestamos activos: {reporte['prestamos_activos']}")
print(f"Multas pendientes: ₡{reporte['multas_pendientes']:,.2f}")
print("="*60)

# Salida:
# ============================================================
# REPORTE DE BIBLIOTECA: Biblioteca Central nablabs
# ============================================================
# Total de libros: 15
# Libros disponibles: 12
# Libros prestados: 3
# Total de usuarios: 8
# Prestamos activos: 3
# Multas pendientes: ₡4,500.00
# ============================================================

# Libros mas prestados
top_libros = biblioteca.libros_mas_prestados(limite=3)
print("\nTop 3 Libros Mas Prestados:")
for i, (libro, cantidad) in enumerate(top_libros, 1):
    print(f"{i}. {libro.titulo} - {cantidad} prestamos")
# Salida:
# Top 3 Libros Mas Prestados:
# 1. 1984 - 12 prestamos
# 2. Sapiens - 8 prestamos
# 3. El Principito - 6 prestamos
\end{lstlisting}

\subsection{Uso de Propiedades y Validación}

\begin{lstlisting}[caption={Ejemplo de propiedades y validación}]
# Crear libro con ISBN invalido (manejo de excepcion)
try:
    libro_invalido = Libro(
        isbn="123",  # ISBN muy corto
        titulo="Libro de Prueba",
        autor="Autor Desconocido",
        editorial="Editorial X",
        año=2024,
        categoria="Prueba"
    )
except ValueError as e:
    print(f"Error: {e}")
# Salida: Error: ISBN invalido: 123

# Usar propiedades calculadas
libro = Libro(
    isbn="978-0747532699",
    titulo="Harry Potter y la Piedra Filosofal",
    autor="J.K. Rowling",
    editorial="Bloomsbury",
    año=1997,
    categoria="Fantasia"
)

# Prestar el libro
usuario = Usuario(
    nombre="Juan Perez",
    identificacion="2-3456-7890",
    email="juan@email.com",
    telefono="7777-7777",
    tipo_usuario="regular"
)

libro.prestar(usuario)

# Acceder a propiedades calculadas
print(f"Libro disponible: {libro.disponible}")  # False (property)
print(f"Dias prestado: {libro.dias_prestado}")  # Calculado automaticamente

# Intentar prestar de nuevo (ya prestado)
if not libro.prestar(usuario):
    print("No se puede prestar: libro ya prestado")
# Salida: No se puede prestar: libro ya prestado
\end{lstlisting}

\subsection{Ejecución Completa del Menú}

\begin{lstlisting}[caption={Ejemplo de interacción con el menú}]
============================================================
Sistema de Gestion de Biblioteca
NLPY-1 - Semana 8: Programacion Orientada a Objetos
============================================================

============================
MENU PRINCIPAL
============================
1. Gestion de libros
2. Gestion de usuarios
3. Realizar prestamo
4. Realizar devolucion
5. Buscar libros
6. Reportes y estadisticas
7. Salir

Seleccione una opcion: 1

============================
GESTION DE LIBROS
============================
1. Agregar libro
2. Buscar libro por ISBN
3. Listar todos los libros
4. Listar libros disponibles
5. Volver

Seleccione una opcion: 1

--- Agregar Nuevo Libro ---
ISBN: 978-0451524935
Titulo: 1984
Autor: George Orwell
Editorial: Signet Classic
Año: 1950
Categoria: Ficcion

¡Libro agregado exitosamente!
ISBN: 978-0451524935
Titulo: '1984' por George Orwell (1950)

Presione Enter para continuar...

============================
MENU PRINCIPAL
============================
...

Seleccione una opcion: 3

--- Realizar Prestamo ---
Identificacion del usuario: 1-2345-6789
ISBN del libro: 978-0451524935

¡Prestamo realizado exitosamente!
Usuario: Maria Rodriguez
Libro: 1984
Fecha de prestamo: 2024-10-24
Fecha limite de devolucion: 2024-11-07
Dias permitidos: 14

Prestamos actuales del usuario: 1/3

============================
MENU PRINCIPAL
============================
...

Seleccione una opcion: 6

============================
REPORTES Y ESTADISTICAS
============================

Biblioteca: Biblioteca Central nablabs
---------------------------
Total de libros: 8
Libros disponibles: 6
Libros prestados: 2
Total de usuarios: 5
Prestamos activos: 2
Total historico de prestamos: 23
Multas pendientes: ₡2,500.00

Top 5 Libros Mas Prestados:
1. 1984 (George Orwell) - 7 prestamos
2. Cien Años de Soledad (Gabriel Garcia Marquez) - 5 prestamos
3. El Principito (Antoine de Saint-Exupery) - 4 prestamos
4. Sapiens (Yuval Noah Harari) - 3 prestamos
5. Don Quijote (Miguel de Cervantes) - 2 prestamos

Presione Enter para continuar...
\end{lstlisting}

\section{Defensa del Código (60\% de la Nota)}

\subsection{Aspectos Técnicos a Dominar}

\textbf{1. Fundamentos de POO (15 puntos)}
\begin{itemize}
    \item ¿Qué es una clase? ¿Qué es un objeto/instancia?
    \item Diferencia entre atributos de clase y de instancia
    \item ¿Para qué sirve \texttt{self}? ¿Por qué es necesario?
    \item ¿Qué hace el método \texttt{\_\_init\_\_}?
    \item Explicar el concepto de encapsulamiento
\end{itemize}

\textbf{2. Herencia y Polimorfismo (15 puntos)}
\begin{itemize}
    \item ¿Qué es herencia? ¿Cuándo usarla?
    \item ¿Para qué sirve \texttt{super()}?
    \item Diferencia entre sobreescritura (override) y sobrecarga
    \item ¿Qué es polimorfismo? Ejemplos en tu código
    \item Herencia simple vs herencia múltiple
\end{itemize}

\textbf{3. Métodos Especiales (10 puntos)}
\begin{itemize}
    \item Diferencia entre \texttt{\_\_str\_\_} y \texttt{\_\_repr\_\_}
    \item ¿Qué hace \texttt{\_\_eq\_\_}? ¿Cuándo se llama?
    \item Otros métodos especiales: \texttt{\_\_len\_\_}, \texttt{\_\_getitem\_\_}
    \item ¿Por qué implementar métodos especiales?
\end{itemize}

\textbf{4. Propiedades y Decoradores (10 puntos)}
\begin{itemize}
    \item ¿Qué es \texttt{@property}? ¿Por qué usarlo?
    \item Diferencia entre property y atributo normal
    \item ¿Cómo implementar un setter?
    \item \texttt{@classmethod} vs \texttt{@staticmethod} vs método de instancia
    \item Ventajas de usar propiedades para validación
\end{itemize}

\textbf{5. Diseño y Buenas Prácticas (10 puntos)}
\begin{itemize}
    \item Principio de responsabilidad única
    \item Composición vs herencia: ¿cuándo usar cada una?
    \item Atributos privados (\_\_) vs protegidos (\_)
    \item Nomenclatura de clases y métodos
    \item Documentación de clases con docstrings
\end{itemize}

\subsection{Preguntas de Defensa Típicas}

\begin{itemize}
    \item "Explica la jerarquía de herencia en tu sistema"
    \item "¿Por qué Usuario hereda de Persona? ¿Qué gana con eso?"
    \item "Muéstrame dónde usas polimorfismo en tu código"
    \item "¿Qué pasa si llamas \texttt{libro.prestar()} cuando ya está prestado?"
    \item "Explica cómo funciona tu property \texttt{dias\_retraso}"
    \item "¿Por qué usaste composición en la clase Prestamo?"
    \item "Modifica tu código para agregar un método que calcule el préstamo promedio"
\end{itemize}

\section{Rúbrica de Evaluación}

\subsection{Funcionalidad del Código (40 puntos)}

\textbf{Implementación de Clases (40 puntos):}
\begin{itemize}
    \item 5 puntos: Clase Persona con atributos y métodos básicos
    \item 8 puntos: Clases Usuario y Bibliotecario con herencia correcta
    \item 8 puntos: Clase Libro con propiedades y validación
    \item 7 puntos: Clase Prestamo con cálculos de multa
    \item 12 puntos: Clase Biblioteca con todas las funcionalidades CRUD
\end{itemize}

\subsection{Defensa del Código (60 puntos)}

\textbf{Fundamentos de POO (15 puntos):}
\begin{itemize}
    \item 5 puntos: Explica clases, objetos y self
    \item 5 puntos: Comprende constructores y atributos
    \item 5 puntos: Entiende encapsulamiento
\end{itemize}

\textbf{Herencia y Polimorfismo (15 puntos):}
\begin{itemize}
    \item 6 puntos: Implementa herencia correctamente
    \item 5 puntos: Usa super() apropiadamente
    \item 4 puntos: Demuestra polimorfismo
\end{itemize}

\textbf{Métodos Especiales (10 puntos):}
\begin{itemize}
    \item 4 puntos: Implementa \_\_str\_\_ y \_\_repr\_\_
    \item 3 puntos: Implementa \_\_eq\_\_ u otros
    \item 3 puntos: Comprende cuándo se llaman
\end{itemize}

\textbf{Propiedades y Decoradores (10 puntos):}
\begin{itemize}
    \item 4 puntos: Usa @property correctamente
    \item 3 puntos: Implementa setters con validación
    \item 3 puntos: Usa @classmethod o @staticmethod
\end{itemize}

\textbf{Diseño y Calidad (10 puntos):}
\begin{itemize}
    \item 3 puntos: Jerarquía coherente y lógica
    \item 3 puntos: Nombres descriptivos y convenciones
    \item 2 puntos: Documentación con docstrings
    \item 2 puntos: Código organizado y limpio
\end{itemize}

\subsection{Penalizaciones}

\begin{itemize}
    \item \textbf{-5 puntos}: No usar herencia cuando es apropiado
    \item \textbf{-5 puntos}: No implementar métodos especiales básicos
    \item \textbf{-4 puntos}: No usar \texttt{super()} en constructores de clases derivadas
    \item \textbf{-3 puntos}: Acceder directamente a atributos privados
    \item \textbf{-3 puntos}: No validar datos en constructores
    \item \textbf{-2 puntos}: Clases sin docstrings
    \item \textbf{-2 puntos}: Nomenclatura incorrecta (clases no en PascalCase)
\end{itemize}

\section{Desafíos Opcionales (Puntos Extra)}

\subsection{Nivel Intermedio (+10 puntos cada uno)}

\textbf{1. Sistema de Categorías de Libros}
\begin{itemize}
    \item Crear clase Categoria con atributos y métodos
    \item Implementar jerarquía de categorías (ficción, no-ficción, etc.)
    \item Búsqueda por categoría jerárquica
\end{itemize}

\textbf{2. Reservas de Libros}
\begin{itemize}
    \item Clase Reserva para cuando libro no está disponible
    \item Lista de espera por libro
    \item Notificación automática cuando se devuelve
\end{itemize}

\textbf{3. Historial Detallado}
\begin{itemize}
    \item Registrar todo el historial de préstamos
    \item Estadísticas por usuario (libros más leídos, categorías favoritas)
    \item Generar reporte de historial
\end{itemize}

\subsection{Nivel Avanzado (+15 puntos cada uno)}

\textbf{4. Sistema de Membresías}
\begin{itemize}
    \item Clases para diferentes tipos de membresía (básica, premium, VIP)
    \item Beneficios por nivel (días de préstamo, cantidad de libros)
    \item Cálculo de cuotas y renovaciones
\end{itemize}

\textbf{5. Persistencia con JSON}
\begin{itemize}
    \item Guardar y cargar biblioteca desde archivos JSON
    \item Serialización de objetos
    \item Métodos de clase para crear desde archivo
\end{itemize}

\textbf{6. Sistema de Recomendaciones}
\begin{itemize}
    \item Recomendar libros basados en historial
    \item Algoritmo simple de similitud
    \item Top 5 recomendaciones personalizadas
\end{itemize}

\subsection{Nivel Experto (+20 puntos)}

\textbf{7. Patrón Observer para Notificaciones}
\begin{itemize}
    \item Implementar patrón Observer
    \item Notificar eventos (préstamo, devolución, multa)
    \item Sistema de suscriptores
\end{itemize}

\textbf{8. Interfaz Gráfica con tkinter}
\begin{itemize}
    \item Crear GUI básica para el sistema
    \item Mantener toda la lógica POO
    \item Ventanas para cada operación principal
\end{itemize}

\section{Recursos y Referencias}

\subsection{Documentación}

\begin{itemize}
    \item Tutorial POO oficial: \url{https://docs.python.org/3/tutorial/classes.html}
    \item Métodos especiales: \url{https://docs.python.org/3/reference/datamodel.html}
    \item Decoradores property: \url{https://docs.python.org/3/library/functions.html#property}
\end{itemize}

\subsection{Consejos para el Desarrollo}

\begin{enumerate}
    \item \textbf{Diseña primero}: Dibuja tu jerarquía de clases antes de programar
    \item \textbf{Empieza simple}: Implementa clases básicas primero, luego agrega complejidad
    \item \textbf{Prueba cada clase}: Verifica cada clase individualmente antes de integrar
    \item \textbf{Usa docstrings}: Documenta cada clase y método importante
    \item \textbf{Piensa en objetos}: Modela entidades del mundo real
    \item \textbf{Valida siempre}: No asumas que los datos son correctos
\end{enumerate}

\section{Criterios de Entrega}

\subsection{Archivos Requeridos}

\begin{enumerate}
    \item \textbf{biblioteca.py}: Código principal con todas las clases
    \item \textbf{pruebas.py}: Script con casos de prueba
    \item \textbf{diagrama\_clases.txt}: Diagrama de jerarquía de clases (texto ASCII)
    \item \textbf{README.txt}: Documentación del proyecto
\end{enumerate}

\subsection{Contenido del README}

\begin{itemize}
    \item Nombre del estudiante
    \item Diagrama de clases (descripción textual)
    \item Instrucciones de uso
    \item Ejemplos de operaciones principales
    \item Desafíos opcionales implementados
    \item Decisiones de diseño importantes
\end{itemize}

\vspace{1cm}

\begin{center}
\textit{"La programación orientada a objetos no es solo sobre sintaxis,}\\
\textit{es sobre pensar en términos de responsabilidades y relaciones."}\\
\vspace{0.3cm}
--- Adaptación de principios de diseño OOP
\end{center}

\vspace{0.5cm}

\begin{center}
\textbf{¡Éxito con tu sistema de biblioteca!}\\
\textit{Recuerda: Un buen diseño orientado a objetos hace el código más mantenible y extensible.}
\end{center}

\end{document}
